\section{معادلات العزوم النهائية للمفاصل}

بعد إهمال قوى كوريوليس والقوى النابذة المركزية، يصبح النموذج الديناميكي المستخدم لحساب عزوم المفاصل بالشكل:

\[
\tau =
\left(D(q)+I_{act}\right)\ddot{q}
+
D_g(q)
\]

حيث تمثل $D(q)$ مصفوفة العطالة، وتمثل $I_{act}$ مصفوفة عزوم القصور الذاتي للمشغلات، بينما يمثل $D_g(q)$ متجه الجاذبية.

وبالتوسع على مستوى كل مفصل من مفاصل الذراع الروبوتية ذات خمس درجات حرية، تكون معادلات العزوم النهائية كما يلي:

\begin{align*}
\tau_1 &=
D_{11}\ddot{\theta}_1
+
I_{act}\ddot{\theta}_1
+
D_{12}\ddot{\theta}_2
+
D_{13}\ddot{\theta}_3
+
D_{14}\ddot{\theta}_4
+
D_{15}\ddot{\theta}_5
+
D_1
\\[0.3cm]
\tau_2 &=
D_{21}\ddot{\theta}_1
+
D_{22}\ddot{\theta}_2
+
I_{act}\ddot{\theta}_2
+
D_{23}\ddot{\theta}_3
+
D_{24}\ddot{\theta}_4
+
D_{25}\ddot{\theta}_5
+
D_2
\\[0.3cm]
\tau_3 &=
D_{31}\ddot{\theta}_1
+
D_{32}\ddot{\theta}_2
+
D_{33}\ddot{\theta}_3
+
I_{act}\ddot{\theta}_3
+
D_{34}\ddot{\theta}_4
+
D_{35}\ddot{\theta}_5
+
D_3
\\[0.3cm]
\tau_4 &=
D_{41}\ddot{\theta}_1
+
D_{42}\ddot{\theta}_2
+
D_{43}\ddot{\theta}_3
+
D_{44}\ddot{\theta}_4
+
I_{act}\ddot{\theta}_4
+
D_{45}\ddot{\theta}_5
+
D_4
\\[0.3cm]
\tau_5 &=
D_{51}\ddot{\theta}_1
+
D_{52}\ddot{\theta}_2
+
D_{53}\ddot{\theta}_3
+
D_{54}\ddot{\theta}_4
+
D_{55}\ddot{\theta}_5
+
I_{act}\ddot{\theta}_5
+
D_5
\end{align*}

حيث تمثل $\tau_i$ العزم المطلوب عند المفصل $i$، وتمثل $D_{ij}$ عناصر مصفوفة العطالة، بينما تمثل $D_i$ مركبة الجاذبية الموافقة للمفصل نفسه.

تمت إضافة الحد $I_{act}$ إلى المعادلات السابقة لأخذ عزم القصور الذاتي للمشغلات بعين الاعتبار، حيث يضاف هذا الحد إلى العنصر القطري الموافق لكل مفصل.

وبذلك تكون معادلات العزوم النهائية قد أصبحت معتمدة على ثلاثة مؤثرات رئيسية فقط: عطالة الوصلات، وقصور المشغلات، وتأثير الجاذبية.

للاطلاع على المعادلات التفصيلية الناتجة عن التعويض في نموذج العزوم.
\begin{center}
\href{file:///D:/Projects/latex/Studying_DIY_Arm/resources/equations/torque-equations.txt}
{معادلات العزوم}
\end{center}








\subsection{تحليل عزوم التثبيت الساكنة عند الوضع الصفري}

بعد الحصول على معادلات العزوم النهائية، تم اختبار النموذج الديناميكي عند الوضع الصفري الساكن، وذلك بهدف تحليل تأثير الجاذبية على مفاصل الذراع الروبوتية. في هذا الاختبار تم تعويض جميع زوايا المفاصل وتسارعاتها الزاوية بالقيم الصفرية:

\[
\theta_1=\theta_2=\theta_3=\theta_4=\theta_5=0
\]

\[
\ddot{\theta}_1=\ddot{\theta}_2=\ddot{\theta}_3=\ddot{\theta}_4=\ddot{\theta}_5=0
\]

وبما أن التسارعات الزاوية معدومة، فإن حدود العطالة تختفي من معادلة العزم:

\[
\left(D(q)+I_{act}\right)\ddot{q}=0
\]

وبالتالي يصبح العزم الناتج عند كل مفصل ناتجًا فقط عن تأثير الجاذبية:

\[
\tau = D_g(q)
\]

أي أن القيم الناتجة تمثل عزوم التثبيت الساكنة اللازمة للمحافظة على الذراع في الوضع الصفري دون حركة.

وقد أعطى النموذج القيم التالية:

\[
\tau(0)=
\begin{bmatrix}
0\\
0.4894\\
0.5572\\
-0.1008\\
0.1393
\end{bmatrix}
\ \mathrm{N.m}
\]

وتتم مناقشة هذه القيم كما يلي:


توضح النتائج أن عزم المفصل الأول معدوم في حالة السكون عند الوضع الصفري. ويعود ذلك إلى أن محور دوران المفصل الأول موازي لاتجاه الجاذبية، وبالتالي فإن وزن الذراع لا يولد عزم قيادة حول هذا المحور. ومع ذلك، قد يسبب وزن الذراع أحمالًا ميكانيكية على بنية القاعدة، مثل قوى رد الفعل أو عزوم الانحناء، لكنها لا تظهر كعزم دوران مطلوب من محرك المفصل الأول.

أما المفصلان الثاني والثالث فقد أعطيا أكبر قيم للعزم الساكن، حيث بلغت قيمة العزم عند المفصل الثاني:

\[
\tau_2 = 0.4894\ \mathrm{N.m}
\]

وبلغت عند المفصل الثالث:

\[
\tau_3 = 0.5572\ \mathrm{N.m}
\]

ويرجع ذلك إلى أن هذين المفصلين يتحملان تأثير وزن معظم الوصلات الواقعة بعدهما. كما أن مراكز كتل هذه الوصلات لا تقع على محاور الدوران الخاصة بهذين المفصلين، مما يؤدي إلى توليد أذرع عزم واضحة بفعل الجاذبية. لذلك تكون عزوم التثبيت المطلوبة عند هذه المفاصل أكبر من عزوم المفاصل القريبة من نهاية الذراع.

بالنسبة للمفصل الرابع، فقد كانت قيمة العزم:

\[
\tau_4 = -0.1008\ \mathrm{N.m}
\]

ويلاحظ أن مقدار هذه القيمة صغير نسبيًا مقارنة بعزمي المفصلين الثاني والثالث:

\[
|\tau_4| = 0.1008\ \mathrm{N.m}
\]

ويرجع ذلك إلى أن المفصل الرابع لا يحمل معظم الذراع، وإنما يتأثر أساسًا بوزن الأجزاء الطرفية الواقعة بعده، مثل الوصلة الأخيرة والقبضة. وبما أن الكتلة المؤثرة بعد هذا المفصل أقل، كما أن مراكز كتل هذه الأجزاء أقرب إلى محور دورانه، فإن ذراع العزم الناتج عن الجاذبية يكون أصغر. لذلك تظهر قيمة العزم عند المفصل الرابع أقل من عزوم المفاصل الأقرب إلى القاعدة.

أما المفصل الخامس فقد أعطى قيمة:

\[
\tau_5 = 0.1393\ \mathrm{N.m}
\]

وهي قيمة صغيرة نسبيًا، لأن هذا المفصل يقع قرب نهاية الذراع ويتأثر بشكل أساسي بوزن الجزء النهائي والقبضة فقط، ولا يحمل عددًا كبيرًا من الوصلات كما هو الحال في المفاصل القريبة من القاعدة.

ومن المهم الإشارة إلى أن الوضع الصفري في هذا الاختبار يعني أن متغيرات المفاصل في نموذج DH تساوي صفرًا:

\[
\theta_i=0
\]

مع بقاء الإزاحات الزاوية الثابتة الموجودة في جدول DH، مثل:

\[
\theta_2+\frac{\pi}{2},
\qquad
\theta_4+\pi,
\qquad
\theta_5+\frac{\pi}{2}
\]